% Context from chapters-tex/variational-priors.tex; for mathematical reference, not compilation.
\section{PDE and Energy Minimization}

Minimizing a smooth prior by gradient descent produces an image-smoothing flow.

                                                                      
\subsection{General Flows}

The gradient of a differentiable prior $J : \RR^N \rightarrow \RR$ is the vector $\grad J(f)$ that describes its local first-order variation:
\eq{
	J(f+\epsilon) = J(f) + \dotp{\epsilon}{\grad J(f)} + o(\norm{\epsilon}).
}

Gradient descent seeks to decrease a smooth discrete energy by iterating
\eql{\label{eq-gradmin-flow}
	f^{(k+1)} = f^{(k)}  - \tau \grad J(f^{(k)}),
}
where the step size $\tau$ is chosen to ensure descent. The quadratic and smoothed TV energies below admit explicit sufficient bounds on $\tau$.

To pass to continuous time, associate iteration $k$ with time $t=k\tau$ and let $\tau$ tend to zero. The limiting flow
\eq{
	t > 0 \longmapsto f_t \in \RR^N
}
satisfies the following differential equation, an ODE for a finite-dimensional image and a PDE when space is continuous:
\eql{\label{eq-pde-flow}
	\pd{f_t}{t} = - \grad J(f_t)
	\qandq
	f_0 = f.
}
Gradient descent is an explicit time discretization of this differential equation at times $t_k = k  \tau$.

\myfigure{
\tabquatre{
\image{variational}{.23}{flow-heat-1}&
\image{variational}{.23}{flow-heat-4}&
\image{variational}{.23}{flow-heat-2}&
\image{variational}{.23}{flow-heat-3}\\
\image{variational}{.23}{flow-tv-4}&
\image{variational}{.23}{flow-tv-3}&
\image{variational}{.23}{flow-tv-2}&
\image{variational}{.23}{flow-tv-1}
}
}{ 
    Heat flow (top) and TV flow (bottom): images $f_t$ at increasing times $t$.
}{fig-flow}

                                                                      
\subsection{Heat Flow}

Heat flow results from applying \eqref{eq-pde-flow} to the Sobolev energy $\Jsob(f)$, defined in \eqref{eq-dfn-sob-energy} for functions and in \eqref{eq-discr-sobolev} for discrete images.

Expanding the quadratic energy gives
\eq{
	J(f+\epsilon) = \frac{1}{2}\norm{\nabla f+\nabla \epsilon}^2 = J(f) - \dotp{\Delta f}{\epsilon} + \frac{1}{2}\norm{\nabla\epsilon}^2,
}
so that 
\eq{
	\grad \Jsob(f) = -\Delta f.
}
Figure \ref{fig-discrete-operators-lapl}, left, shows an image Laplacian. Its magnitude is typically large near edges, with either sign.

The heat flow is thus 
\eql{\label{eq-heat-flow-cont}
	\pd{f_t}{t}(x) = - (\grad J(f_t))(x) = \De f_t(x)
	\qandq
	f_0 = f.
}

\myfigure{
\tabdeux{
\image{variational}{.3}{discrete-operators-lapl}&
\image{variational}{.3}{discrete-operators-lapltv}\\
$\diverg(\nabla f)$ & $\diverg(\nabla f / \norm{\nabla f})$ 
}
}{ 
	Discrete Laplacian (left) and negative discrete TV gradient (right).
}{fig-discrete-operators-lapl}

  
\paragraph{Continuous in space.}

For continuous images on the unbounded domain $\RR^2$, the PDE \eqref{eq-heat-flow-cont} has an explicit solution: convolution with a Gaussian kernel whose variance grows with time
\eql{\label{eq-heat-gaussian}
	f_t = f \star h_t \qwhereq 
	h_t(x) = \frac{1}{4 \pi t} e^{-\frac{\norm{x}^2}{4t}}.
}	
Convolution with this widening kernel progressively smooths the image.

  
\paragraph{Discrete in space.}

The gradient flow of the discrete Sobolev energy \eqref{eq-discr-sobolev} is
\eq{
	\pd{f_{n,t}}{t} = - (\grad J(f_t))_n = (\De f_t)_n.
}
Discretizing time as in \eqref{eq-gradmin-flow} gives the fully discrete iteration
\eq{	
	f^{(k+1)}_n = f^{(k)}_n + \tau \Big( \sum_{p \in V_4(n)} f_p^{(k)} - 4 f_n^{(k)} \Big) = (f^{(k)} \star h)_n
}
where $V_4(n)$ denotes the four neighbors of pixel $n$. Each iteration convolves the current image with the same filter:
\eq{
	f^{(k)} = f \star h \star \ldots \star h  = f \star^k h.
}
The filter has coefficients $h_0=1-4\tau$ and $h_{\pm e_1}=h_{\pm e_2}=\tau$; all other entries vanish.

This iteration is stable and converges if $0<\tau<1/4$.



                                                                      
\subsection{Total Variation Flows}

  
\paragraph{Total variation gradient.}

The total variation energy $\Jtv$ is nonsmooth, both in its continuous form \eqref{eq-dfn-tv-energy} and in its discrete form \eqref{eq-discr-tv}. In the discrete case, $\Jtv$ is nondifferentiable at any image $f$ with a pixel $n$ where $\nabla f_n = 0$.

If $\nabla f_n\neq0$ at every pixel, the gradient is
\eq{
	(\grad J(f))_n = -\diverg\pa{ \frac{ \nabla f}{ \norm{\nabla f} } }_n
}
The displayed quotient is undefined where the spatial gradient vanishes. Figure \ref{fig-discrete-operators-lapl}, right, shows the negative TV gradient; normalization by small gradient magnitudes makes it sensitive to perturbations in smooth regions.

Nondifferentiability prevents the use of classical gradient descent at such im